
\documentclass[../../../physics-standard_questions_all.tex]{subfiles}

\begin{document}

真空中の誘電率を$\varepsilon_0$とする．

\begin{enumerate}

    \item クーロンの法則の比例定数$k_0$を，$\varepsilon_0$を用いて表せ．
    
    \item 充分長い直線上に一様に電荷が分布しており，
    単位長さあたりの電荷（電荷の線密度）は$\lambda \, (\lambda>0)$である．
    この場合，直線に充分近い点では，電気力線は直線に垂直に一様に出ていると見なせる．
    その条件下で，直線から距離$r$の点での電場の強さ$E(r)$を求めよ．

    \item 充分広い平面上に一様に電荷が分布しており，
    単位面積あたりの電荷（電荷の面密度）は$\sigma \, (\sigma>0)$である．
    この場合，平面に充分近い点では，電気力線は平面に垂直に一様に出ていると見なせる．
    その条件下で，周囲での電場の強さ$E$を求めよ．

    \item 真空中に半径$a$の球殻（厚み無視）があり，一様に帯電している．
    その電荷の総量は$Q \, (Q>0)$である．
    球殻の中心から$r$離れた点での電場の強さ$E(r)$を，
    $0<r<a$の領域と$r>a$の領域に分けて，それぞれ求めよ．

    \item 真空中に半径$a$の荷電球があり，一様に帯電している．
    その電荷の総量は$Q \, (Q>0)$である．
    球殻の中心から$r$離れた点での電場の強さ$E(r)$を，
    $0<r<a$の領域と$r>a$の領域に分けて，それぞれ求めよ．

    \item ともに$+Q$の電荷を持つ2つの点電荷が，ある距離を隔てて固定されている．
    この周囲の電気力線の概略を図示せよ．

\end{enumerate}

\end{document}


